% -- new commands by MiG---

\newcommand{\eg}{e.g.,\xspace}
\newcommand{\Eg}{E.g.,\xspace}
\newcommand{\ie}{i.e.,\xspace}
\newcommand{\Ie}{I.e.,\xspace}
\newcommand{\incl}{incl.\xspace}
\newcommand{\Incl}{Incl.\xspace}

\newcommand{\bspw}{bspw.\xspace}
\newcommand{\Bspw}{Bspw.\xspace}
\newcommand{\bzw}{bzw.\xspace}
\newcommand{\ca}{ca.\xspace}
\newcommand{\Ca}{Ca.\xspace}
\newcommand{\dhe}{d.\,h.\xspace}
\newcommand{\Dhe}{D.\,h.\xspace}
\newcommand{\etc}{etc.\xspace}
\newcommand{\eV}{e.\,V.\xspace}
\newcommand{\evtl}{evtl.\xspace}
\newcommand{\Evtl}{Evtl.\xspace}
\newcommand{\ggf}{ggf.\xspace}
\newcommand{\ia}{i.\,a.\xspace}
\newcommand{\iA}{i.\,A.\xspace}
\newcommand{\IA}{I.\,A.\xspace}
\newcommand{\insbes}{insbes.\xspace}
\newcommand{\Insbes}{Insbes.\xspace}
\newcommand{\inkl}{inkl.\xspace}
\newcommand{\Inkl}{Inkl.\xspace}
\newcommand{\oae}{o.\,ä.\xspace}
\newcommand{\og}{o.\,g.\xspace}
\newcommand{\sog}{sog.\xspace}
\newcommand{\Sog}{Sog.\xspace}
\newcommand{\teilw}{teilw.\xspace}
\newcommand{\Teilw}{Teilw.\xspace}
\newcommand{\ua}{u.\,a.\xspace}
\newcommand{\Ua}{U.\,a.\xspace}
\newcommand{\usw}{usw.\xspace}
\newcommand{\uU}{u.\,U.\xspace}
\newcommand{\UU}{U.\,U.\xspace}
\newcommand{\vgl}{vgl.\xspace}
\newcommand{\Vgl}{Vgl.\xspace}
\newcommand{\zB}{z.\,B.\xspace}
\newcommand{\zb}{z.\,B.\xspace}
\newcommand{\ZB}{Z.\,B.\xspace}
\newcommand{\Zb}{Z.\,B.\xspace}


%Macro for code inclusions
\newcommand{\code}[1]{\lstinline|#1|}

% macros for emphasizing text
\newcommand{\myemph}[1]{\emph{#1}}
\newcommand{\mydef}[1]{\textbf{#1}}
\newcommand{\myexcl}[1]{\textbf{#1}}

% Short command for backslash in text mode:
\newcommand{\bs}{\textbackslash}
% Short command for printing latex commands (appends backslash in front)
\newcommand{\lc}[1]{{\ttfamily\textbackslash #1}}
% Formatting keywords, menu settings, parameters
\newcommand{\printkeyword}[1]{\enquote{\ttfamily #1}}
% Formatting file names or file paths
\newcommand{\printfilepath}[1]{{\ttfamily #1}}
\newcommand{\printfilepathquoted}[1]{\enquote{\ttfamily #1}}
% Formatting name of a software (in bold)
\newcommand{\printswname}[1]{{\bfseries #1}}
% Short command for referencing a latex package in the text (adds an index)
\newcommand{\pkg}[1]{{\ttfamily\index{#1}\index{Paket!#1}#1}}

%\newcommand{\bild}[6]% Bild-Pfadname, Beschriftung, Label, Breite, Kurzbeschriftung (für Abbildungsverzeichnis), optional: Platzierung
%{%
%  \begin{figure}[#6]%
%		\Centering%
%        \includegraphics*[width=#4]{#1}%
%        \caption[#5]{\label{#3} #2}%
% \end{figure}
%}



% 1.) Seitliche Kommentare/Abschnitts-Untertitel
%   a.) Für Verwendung im Text (Float-Variante)
\newcommand{\floatmarginnote}[1]{%
\ifthenelse{\boolean{showMarginNotes}}{%if margin notes should be displayed:
%Variante mit Kapitälchen
%\marginpar[\flushleft{\textcolor{gray}{\textsc{#1}}}]{\flushleft{\textcolor{gray}{\textsc{#1}}}}%
%Variante ohne Kapitälchen
\marginpar[\flushleft{\textcolor{gray}{#1}}]{\flushleft{\textcolor{gray}{#1}}}%
}{%else
\relax%
}}

%   b.) Non-Float-Variante (für Verwendung in Gleichungen):
\newcommand{\nonfloatmarginnote}[1]{%
\ifthenelse{\boolean{showMarginNotes}}{%if margin notes should be displayed:
\marginnote[\RaggedRight#1]{\RaggedRight#1}%
%%Variante mit Kapitälchen (obsolet, s.u.)
%\marginnote[\textcolor{gray}{\textsc{#1}}]{\textcolor{gray}{\textsc{#1}}}
%%Variante mit Kapitälchen wird direkt durch die folgenden Befehle in der Präambel gesetzt:
%%\renewcommand*{\raggedleftmarginnote}{}
%%\renewcommand*{\marginfont}{\color{gray}\sffamily\scshape}
}{%else
\relax%
}}

\newlength{\marginwidth}
\setlength{\marginwidth}{\marginparwidth}
\addtolength{\marginwidth}{\marginparsep}


%Breite der Grafiken in einer fbox:
\newcommand{\linewidthwithoutfbox}{\linewidth-2\fboxsep-2\fboxrule}


%%%%%%%%%%%%%%%%%%%%%%%%%%%%%%%%%%%%%%%%%%%%%%%%%%%%%%%%%%%%%%%%%%
%% Notationstabelle mit zwei Spalten mit variablen Spaltenbreite %
%%%%%%%%%%%%%%%%%%%%%%%%%%%%%%%%%%%%%%%%%%%%%%%%%%%%%%%%%%%%%%%%%%
%% Die Breite der ersten Spalte ist als optionaler Argument der Umgebung angebbar.
%% Default ist 4em.
%% Die Breite der zweiten Spalte wird automatisch berechnet.
\newenvironment{myNotationTable}[1][4em]{%
		\begingroup%
		\setlength{\tabcolsep}{0pt}% Kein Einzug bei Notation
		\renewcommand{\arraystretch}{1.3}% Etwas mehr Abstand zwuischen den Zeilen, damit sie nicht zusammenfallen
		%% Forderung des KSP-Verlages: keine Worttrennung in einer Tabelle, daher \raggedright statt \RaggedRight
		\begin{longtable}{>{\raggedright\arraybackslash}p{\dimexpr#1-2\tabcolsep\relax}>{\raggedright\arraybackslash}p{\dimexpr\linewidth-#1-2\tabcolsep\relax}}%
	}{%
		\end{longtable}%
		\endgroup%
	}

%%%%%%%%%%%%%%%%%%%%%%%%%%%%%%%%%%%%%%%%%%%%%%%%%%%%%%%%%%%%%%%%%%
%% Notationstabelle mit drei Spalten mit variablen Spaltenbreite %
%%%%%%%%%%%%%%%%%%%%%%%%%%%%%%%%%%%%%%%%%%%%%%%%%%%%%%%%%%%%%%%%%%
%% Die Breite der ersten Spalte ist als optionaler Argument der Umgebung angebbar.
%% Default ist 30% der Seitenbreite (0.3\linewidth).
%% Die Breite der zweiten Spalte wird automatisch berechnet.
%% Die Breite der letzten Spalte ist fest auf 10% der Seitenbreite gesetzt (0.1\linewidth).
\newenvironment{myNotationDescTable}[1][0.3\linewidth]{%
		\begingroup%
		% Breite der dritten Spalte (Formelspalte) (anpassbar, 
		% aktuell 10% der Seitenbreite (0.1\linewidth)):
		\newlength{\formulaseplength}
		\setlength{\formulaseplength}{0.1\linewidth}
		\setlength{\tabcolsep}{0pt}% Kein Einzug bei Notation
		\renewcommand{\arraystretch}{1.3}%
		%% Forderung des KSP-Verlages: keine Zeilenumbrüche, daher \raggedright statt \RaggedRight
		\begin{longtable}{>{\raggedright\arraybackslash}p{\dimexpr#1-2\tabcolsep}>{\raggedright\arraybackslash}p{\dimexpr\linewidth-\formulaseplength-#1-2\tabcolsep\relax}>{\arraybackslash}p{\formulaseplength-2\tabcolsep}}%
	}{%
		\end{longtable}%
		\endgroup%
	}

%%%%%%%%%%%%%%%%%%%%%%%%%%%%%%%%%%%%%%%%%%%%%%%
%% Zeilen der dreispaltigen Notationstabelle: %
%%%%%%%%%%%%%%%%%%%%%%%%%%%%%%%%%%%%%%%%%%%%%%%
%% 3 Zellen, dritte Zelle ist in Mathe-Format
\newcommand{\myNotationDescTableEntry}[3]{{#1}&{#2}&{$#3$}\\}
%% Alle 3 Zellen im Textformat
\newcommand{\myNotationDescTableEntryText}[3]{{#1}&{#2}&{#3}\\}

%%%%%%%%%%%%%%%%%%%%%%%%%%%%%%%%%%%%%%%%%%%%%%%
%% Zeilen der zweispaltigen Notationstabelle: %
%%%%%%%%%%%%%%%%%%%%%%%%%%%%%%%%%%%%%%%%%%%%%%%
%% 2 Zellen, beide Zellen im Text-Fomrat
\newcommand{\myNotationTableEntryText}[2]{{#1}&{#2}\\}
%% 2 Zellen, erste Zelle ist in Mathe-Format
\newcommand{\myNotationTableEntryMath}[2]{{$#1$}&{#2}\\}


\makeatletter
% Record only entries actually rendered in a contents/list file. The label
% also lets headings outside the TOC remain plain text on subsequent runs.
\let\oldcontentsline\contentsline
\renewcommand{\contentsline}[4]{%
	\oldcontentsline{#1}{%
		% Multiple TODO entries can share one destination.
		\ifcsundef{my@tocseen:#4}{%
			\csgdef{my@tocseen:#4}{}%
			\Hy@raisedlink{\hypertarget{toc.#4}{}}%
			\begingroup
				\edef\@currentHref{toc.#4}%
				\label{toc-entry:#4}%
			\endgroup
		}{}%
		#2}{#3}{#4}%
}

\newcommand{\my@tocbacklink}[2]{%
	\ifthenelse{\boolean{showTableOfContents}}{%
		\@ifundefined{r@toc-entry:#1}{#2}{\hyperlink{toc.#1}{#2}}%
	}{#2}%
}

% Glossaries, the index and the TODO list add their TOC entry before their
% starred heading; front-matter chapters add it afterwards. Remember the
% former's destination rather than using the new anchor made by chapter*.
\let\my@lastchapterhref\@empty
\let\my@pendingchapterhref\@empty
\let\my@oldaddcontentsline\addcontentsline
\renewcommand{\addcontentsline}[3]{%
	\my@oldaddcontentsline{#1}{#2}{#3}%
	\ifstrequal{#1}{toc}{%
		\begingroup
			\edef\my@toclevel{#2}%
			\ifdefstring{\my@toclevel}{chapter}{%
				\ifx\@currentHref\my@lastchapterhref\else
					\global\let\my@pendingchapterhref\@currentHref
				\fi
			}{}%
		\endgroup
	}{}%
}

\newcommand{\my@chaptertitle}[1]{%
	\let\my@backhref\@currentHref
	\ifx\my@pendingchapterhref\@empty\else
		\let\my@backhref\my@pendingchapterhref
	\fi
	\global\let\my@pendingchapterhref\@empty
	\global\let\my@lastchapterhref\@currentHref
	\my@tocbacklink{\my@backhref}{#1}%
}

% KOMA's formatting hooks run after Hyperref assigns the heading anchor.
% Keep the sectioning commands intact, including their optional arguments.
\let\my@oldchapterlinesformat\chapterlinesformat
\renewcommand{\chapterlinesformat}[3]{%
	\my@oldchapterlinesformat{#1}{#2}{\my@chaptertitle{#3}}%
}
\let\my@oldchapterlineswithprefixformat\chapterlineswithprefixformat
\renewcommand{\chapterlineswithprefixformat}[3]{%
	\my@oldchapterlineswithprefixformat{#1}{#2}{\my@chaptertitle{#3}}%
}
\let\my@oldsectionlinesformat\sectionlinesformat
\renewcommand{\sectionlinesformat}[4]{%
	\my@oldsectionlinesformat{#1}{#2}{#3}{\my@tocbacklink{\@currentHref}{#4}}%
}
\let\my@oldsectioncatchphraseformat\sectioncatchphraseformat
\renewcommand{\sectioncatchphraseformat}[4]{%
	\my@oldsectioncatchphraseformat{#1}{#2}{#3}{\my@tocbacklink{\@currentHref}{#4}}%
}
\makeatother


\defbibenvironment{thesislist}
{\list
	{\printtext[labelnumberwidth]{%
			\printfield{labelprefix}%
			\printfield{labelnumber}}}
	{\setlength{\labelwidth}{\labelnumberwidth}%
		\setlength{\leftmargin}{\labelwidth}%
		\setlength{\labelsep}{\biblabelsep}%
		\addtolength{\leftmargin}{\labelsep}%
		\setlength{\itemsep}{\bibitemsep}%
		\setlength{\parsep}{\bibparsep}}%
	\renewcommand*{\makelabel}[1]{\hss##1}}
{\endlist}
{\item}

\DefineBibliographyStrings{english}{
	mathesis = {Master's Thesis},
	bathesis = {Bachelor's Thesis},
}

\DeclareCiteCommand{\citePrefixed}
{\usebibmacro{prenote}}
{\printtext[bibhyperref]{%
		\mkbibbrackets{%
			\printfield{labelprefix}%
			\printfield{labelnumber}}}}
{\multicitedelim}
{\usebibmacro{postnote}}

% Automatically trimming/cropping PDFs
\newcommand{\includeCroppedPdf}[2][]{%
	%\IfFileExists{./#2-crop.pdf}{}{%hash sum check required to check for pdf-updates
		\immediate\write18{pdfcrop --hires #2 #2-crop.pdf}%}%requires build-option: '--shell-escape'
	\includegraphics[#1]{#2-crop.pdf}}

\newtheoremstyle{research}
{3pt}{3pt}{\itshape}{}{\bfseries}{\\}{.5em}{}

\theoremstyle{research}
\newtheorem{rqu}{Research Question}

\newtcolorbox{cbox_withtitle}[2][]{
	fonttitle=\bfseries,
	title=#2,
	colback = white,
	colframe = black,
	boxrule = 1pt,
	leftrule = 1pt, % left rule weight
	#1
}
\newtcolorbox{cbox_notitle}{
	colback = white,
	colframe = black,
	boxrule = 1pt,
	leftrule = 5pt, % left rule weight
}

\newtcolorbox{rqu_box}[1][]{%
	blend before title=colon hang,
	colback = white,
	colframe = black,
	boxrule = 1pt,
	leftrule = 5pt,
	%left=2pt,
	%right=2pt,
	%boxsep=2pt, %inner padding
	%every float=\centering,
	halign=flush left,
	halign lower=flush right,
	righthand width=0.15\textwidth,
	lower separated=false,
	#1}

% Reference: https://tex.stackexchange.com/questions/340937/tcolorbox-and-figures-adding-description
%\newtcolorbox[blend into=figures]{myfigure_orig}[2][]{%
%	float=htb,
%	blend before title=colon hang,
%	sharp corners,
%	colback=white,
%	colbacktitle=blue!30!black,
%	title={#2},
%	every float=\centering,
%	halign=flush center,
%	halign lower=flush center,
%	lower separated=false,
%	#1}

\newlength{\capheight}
\settoheight{\capheight}{X}

%\newcommand{\mythinarrowup}{\tikz[baseline=-.9ex] \draw[-stealth, thin] (0,-0.115) -- (0,.1);}
%\newcommand{\mythinarrowdown}{\tikz[baseline=-.9ex] \draw[-stealth, thin] (0,0.1) -- (0,-.115);}

\newcommand{\mythinarrowup}{%
	\settoheight{\capheight}{X}%
	\tikz[baseline=-0.ex]
	\draw[-stealth, thin] (0,0) -- (0,\capheight);%
}
\newcommand{\mythinarrowdown}{%
	\settoheight{\capheight}{X}%
	\tikz[baseline=-0.ex]
	\draw[-stealth, thin] (0,\capheight) -- (0,0);%
}

% For ranks in tables or elsewhere:
\newcommand{\ranknum}[1]{\textit{\textbf{\small #1}}}
\newcommand{\rankstar}[2]{\textbf{\small #1\textsuperscript{*}}\,\scriptsize{#2}}
\newcommand{\rankdiff}[2]{\textit{\textbf{\small #1}}\,\scriptsize{#2}}
